\documentclass[11pt,a4paper]{article}

\usepackage{iftex}
\RequirePDFTeX
\usepackage{cmap}
\usepackage[T2A]{fontenc}
\usepackage[utf8]{inputenc}
\usepackage[english,russian]{babel}
\usepackage{PTSerif}
\usepackage{amsmath,amssymb,microtype}
\usepackage[a4paper,left=23mm,right=23mm,top=20mm,bottom=21mm]{geometry}
\usepackage{xcolor,enumitem,fancyhdr,hyperref}

\definecolor{accent}{HTML}{187AA5}
\definecolor{ink}{HTML}{302B28}
\definecolor{muted}{HTML}{6D6560}
\definecolor{tint}{HTML}{EDF4F6}
\definecolor{rulecolor}{HTML}{A9CADA}
\hypersetup{unicode=true,colorlinks=true,urlcolor=accent,linkcolor=accent,
  pdftitle={Домашнее задание 4. Околокруговое приближение. Одноимпульсные манёвры. Базис-вектор.},
  pdfauthor={Чиняев Владимир Николаевич}}
\urlstyle{same}
\color{ink}
\setlength{\parindent}{0pt}
\setlength{\parskip}{3pt}
\linespread{1.08}
\setlist[enumerate]{leftmargin=*,label=\textcolor{accent}{\arabic*.},
  itemsep=3pt,topsep=4pt,parsep=0pt}
\setlist[itemize]{leftmargin=*,itemsep=3pt,topsep=3pt,parsep=0pt}
\setlength{\emergencystretch}{1em}
\pagestyle{fancy}
\fancyhf{}
\renewcommand{\headrulewidth}{0pt}
\renewcommand{\footrulewidth}{0.4pt}
\renewcommand{\footrule}{\hbox to\headwidth{\color{rulecolor}\leaders\hrule height\footrulewidth\hfill}}
\fancyfoot[L]{\footnotesize\color{muted}Введение в маневрирование КА\quad /\quad ДЗ 4}
\fancyfoot[R]{\footnotesize\color{muted}\thepage}

\newcommand{\taskheading}[2]{%
  \par\vspace{5pt}{\large\bfseries\color{accent}Задание #1. #2}\par\smallskip}
\newcommand{\hint}[1]{%
  \par\nobreak\smallskip\noindent
  \begingroup
  \setlength{\fboxsep}{7pt}%
  \setlength{\fboxrule}{0.4pt}%
  \fcolorbox{rulecolor}{tint}{%
    \begin{minipage}{\dimexpr\linewidth-2\fboxsep-2\fboxrule\relax}
    \small\textbf{\color{accent}Подсказка.} #1\par
    \end{minipage}}%
  \endgroup\par\smallskip}

\begin{document}

{\small\color{accent}ВВЕДЕНИЕ В МАНЕВРИРОВАНИЕ КА}\par
\vspace{3pt}
{\huge Домашнее задание 4}\par
{\large Околокруговое приближение.\\Одноимпульсные манёвры. Базис-вектор.}\par
{\small\color{muted}К лекции 5}
\par{\small\bfseries\color{accent}Срок сдачи: 26 октября 2026 года, 23:59}
\vspace{4pt}
{\color{rulecolor}\hrule}
\vspace{5pt}

Решите задачу перехода между близкими компланарными орбитами
в околокруговом приближении.
Задания 1--3 обязательны, задание 4* --- по желанию.
Решение оформите в предоставленном Jupyter Notebook.

\taskheading{1}{Чтение методички}

В методичке \texttt{near-circular-transfer.pdf} прочитайте разделы
1.5 и 1.9, а также 2.2--2.4: импульсное изменение параметров,
постановку задачи перехода, базис-вектор и условия оптимальности.
Раздел 2.1 используйте для справки по обозначениям.

Дайте краткие обоснованные ответы:
\begin{enumerate}
\item Почему выполнение граничных условий само по себе не доказывает
оптимальность манёвра?
\item Какие условия на манёвр и один общий базис-вектор достаточны
для глобального минимума характеристической скорости? Почему проверки модуля только в точках импульсов
недостаточно?
\end{enumerate}

\begingroup
\setlength{\fboxsep}{9pt}
\colorbox{tint}{\begin{minipage}{\dimexpr\linewidth-2\fboxsep\relax}
\textbf{Общие исходные данные}
\[
 \mu=398600{,}44158\ \text{км}^3/\text{с}^2,\qquad
 r_0=a_i=7000\ \text{км},\qquad V_0=\sqrt{\frac{\mu}{r_0}}.
\]
\[
 e_{xi}=e_{yi}=0,\qquad a_f=7007\ \text{км},\qquad
 e_{xf}=0{,}0012,\quad e_{yf}=0{,}0016.
\]
Ось отсчёта углов фиксирована в общей плоскости орбит;
$e_x,e_y$ обеих орбит заданы относительно этой оси.
\[
 \Delta a=\frac{a_f-a_i}{r_0},\qquad
 \Delta e_x=e_{xf}-e_{xi},\qquad \Delta e_y=e_{yf}-e_{yi},
 \qquad \Delta e=\sqrt{(\Delta e_x)^2+(\Delta e_y)^2}.
\]
Компоненты импульса нормируются на $V_0$:
$\Delta\widehat V_r=\Delta V_r/V_0$,
$\Delta\widehat V_t=\Delta V_t/V_0$; нормальная компонента равна нулю.
\end{minipage}}
\endgroup

\vspace{4pt}
\textbf{Работа с ноутбуком и самопроверка}

В \texttt{hw\_4.ipynb} готовы константы, вывод таблиц, график и чекеры
\texttt{check\_task2}, \texttt{check\_task3}, \texttt{check\_star}.
Заполните отмеченные переменные и текстовые ячейки своими решениями.
Чекеры предназначены только для самопроверки. Использовать их не требуется.

При использовании чекеров заданий 2 и 3 передавайте углы в радианах и массивы
нормированных компонент: каждая строка содержит
$(\Delta\widehat V_r,\Delta\widehat V_t)$.
В задании 2 строки относятся к отдельным одноимпульсным решениям,
в задании 3 --- к двум импульсам одного перехода.
Множители передаются в порядке $(\lambda_1,\lambda_2,\lambda_3)$.
Расчёты выполняйте без промежуточного округления.

Чекеры проверяют только численные ответы; выводы формул и доказательства
оформите самостоятельно.

\newpage
{\small\color{muted}Домашнее задание 4\hfill К лекции 5}\par
\vspace{5pt}

\taskheading{2}{Одноимпульсный переход}

Используйте три граничных условия околокругового приближения:
\[
\begin{aligned}
 2\Delta\widehat V_t&=\Delta a,\\
 \Delta\widehat V_r\sin\varphi+2\Delta\widehat V_t\cos\varphi&=\Delta e_x,\\
 -\Delta\widehat V_r\cos\varphi+2\Delta\widehat V_t\sin\varphi&=\Delta e_y.
\end{aligned}
\]
\begin{enumerate}
\item Исключите компоненты импульса и получите уравнение для угла $\varphi$.
Выведите условие существования решения; укажите, когда имеются две точки,
одна точка или нет точек перехода. Вырожденный случай
$\Delta a=\Delta e=0$ здесь не рассматривайте.
\item Для общих исходных данных найдите все углы $0\leq\varphi<2\pi$,
обе компоненты и модуль каждого импульса. Представьте таблицу:
угол в градусах, компоненты и модуль в м/с.
Объясните совпадение модулей и связь найденных углов с пересечением орбит.
\item Сохранив $e_{xf},e_{yf}$, определите число решений при
$a_f=7014$~км и $a_f=7021$~км. Полный расчёт импульсов повторять не нужно.
Классифицируйте эти случаи в околокруговом приближении.
\end{enumerate}
\hint{Умножьте вторую строку на $\cos\varphi$, третью --- на
$\sin\varphi$ и сложите: радиальная компонента исключится.
Для восстановления этой компоненты используйте другую линейную комбинацию.
Места импульсов находите аналитически из уравнений околокругового приближения.}

\taskheading{3}{Два импульса и проверка оптимальности}

Вернитесь к $a_f=7007$~км и остальным общим данным.
Рассмотрите два чисто трансверсальных импульса, приложенных при углах
\[
 \cos\varphi_1=\frac{\Delta e_x}{\Delta e},\qquad
 \sin\varphi_1=\frac{\Delta e_y}{\Delta e},\qquad
 \varphi_2=\varphi_1+\pi\pmod{2\pi}.
\]
\begin{enumerate}
\item Из граничных условий найдите $\Delta\widehat V_{t1}$ и
$\Delta\widehat V_{t2}$, сохраняя знаки компонент.
Проверьте все три условия на $\Delta a,\Delta e_x,\Delta e_y$.
\item Найдите суммарную характеристическую скорость в м/с и сравните
с одноимпульсным переходом.
\item Постройте базис-вектор, положив
$\lambda_1=\lambda_4=\lambda_5=\lambda_6=0$.
Найдите $\lambda_2,\lambda_3$ из условия совпадения
$\boldsymbol s(\varphi_1)$ с единичным вектором первого импульса.
Проверьте это условие и для второго импульса, учитывая его знак.
\item Постройте график $\|\boldsymbol s(\varphi)\|$ при
$0\leq\varphi\leq2\pi$, добавьте линию единичного модуля и отметьте углы приложения импульсов.
Аналитически докажите $\|\boldsymbol s(\varphi)\|\leq1$ на всём обороте.
По достаточным условиям обоснуйте глобальную оптимальность среди манёвров
с произвольным числом импульсов.
\end{enumerate}

Для выбранной постановки используйте
\[
\begin{aligned}
 s_r&=\lambda_3\sin\varphi-\lambda_2\cos\varphi,\\
 s_t&=2\lambda_1+2\lambda_3\cos\varphi+2\lambda_2\sin\varphi,
 \qquad s_z=0.
\end{aligned}
\]
\hint{Для найденных множителей удобно выразить компоненты через
$\varphi-\varphi_1$ и вычислить квадрат модуля.
График иллюстрирует результат; аналитическое неравенство обеспечивает проверку
всех углов между узлами сетки.}

\newpage
{\small\color{muted}Домашнее задание 4\hfill К лекции 5}\par
\vspace{5pt}

\taskheading{4*}{Когда один импульс оптимален?}
\textit{Необязательное задание.}

Решите задачу перехода между компланарными орбитами в околокруговом
приближении для произвольных малых разностей $\Delta a,\Delta e_x,\Delta e_y$.
Найдите оптимальное значение суммарной характеристической скорости.
\[
 J_{\min}(\Delta a,\Delta e_x,\Delta e_y),\qquad
 J=V_0\sum_j\sqrt{(\Delta\widehat V_{rj})^2+(\Delta\widehat V_{tj})^2}.
\]
\begin{enumerate}
\item Обобщите двухимпульсную схему задания 3.
При $\Delta e>0$ используйте направление
$(\Delta e_x,\Delta e_y)$ и противоположную точку орбиты.
Найдите компоненты и суммарную характеристическую скорость для произвольного $\Delta a$.
Отдельно рассмотрите $\Delta e=0$ и полностью нулевые разности.
\item Получите выражение для $J_{\min}$ и докажите его через базис-вектор.
Разберите разные соотношения $|\Delta a|$ и $\Delta e$;
при необходимости используйте разные наборы множителей.
\item Установите, когда существует \emph{ненулевой одноимпульсный переход},
достигающий этого минимума. Отличайте отсутствие одноимпульсного решения
от его существования с большей характеристической скоростью.
При полностью нулевых разностях минимум достигается без манёвров;
этот случай не называйте ненулевым одноимпульсным переходом.
\item Для самопроверки при желании реализуйте в ноутбуке две функции
с аргументами \texttt{da, dex, dey}.
Функция \texttt{j\_min\_star} возвращает $J_{\min}$ в м/с,
а \texttt{one\_impulse\_optimal\_star} --- логическое значение
по критерию предыдущего пункта. Их можно проверить с помощью \texttt{check\_star}.
\end{enumerate}
\hint{Получение допустимой схемы даёт верхнюю оценку характеристической скорости.
Для доказательства минимума найдите базис-вектор, который даёт ту же
нижнюю оценку. Не забывайте брать модули отдельных импульсов;
направление вектора эксцентриситета при $\Delta e=0$ не определено.}

\texttt{check\_star} проверяет обе функции на наборах с разными знаками
$\Delta a$, ориентациями вектора эксцентриситета, граничными и нулевыми случаями.
Совпадение численных ответов на этих примерах не заменяет общего доказательства.

\vspace{5pt}
{\color{rulecolor}\hrule}
\vspace{5pt}
\textbf{Что представить}

Сдайте один ноутбук с ответами задания 1, полными решениями заданий 2 и 3
и, по желанию, задания 4*: выводами формул, таблицами импульсов,
графиком модуля базис-вектора и обоснованием оптимальности.
Использование чекеров и сохранение результатов их вызова не требуются.

Формулы и рассуждения оформите в текстовых ячейках, расчёты и рисунки ---
в кодовых. Перед отправкой перезапустите ядро, выполните ячейки решения по порядку
и сохраните ноутбук с результатами.

Название файла: \texttt{hw\_4\_name\_surname.ipynb}.
Вместо \texttt{name\_surname} укажите свои имя и фамилию латиницей.

\end{document}
